\scanpage{4}
\uline{Lemme 1} Soit $k$ un corps, et supposons que les conditions suivantes soient \unsure{vérifiées|vérifiée}

\deleted{\unsure{1)}} On a $H^2(k,\unsure{\mathscr{K}|\mathcal{K}})=0$ dans les cas suivants \inserted{($\mathscr{K}$ commutatif)}

\begin{enumerate}[label=\alph*)]
\item $\mathscr{K}$ un tore
\item $\mathscr{K}$ \deleted{un \unsure{schéma}} fini ét. sur $k$
\item $\mathscr{K}$ une V.A.
\item $\mathscr{K}$ groupe unipotent connexe annulé par $p$.
\end{enumerate}

Alors pour tte \deleted{\unsure{donnée}} \inserted{suite exacte} de \inserted{\unsure{schémas en}} groupes Alg.
\diagram{\begin{tikzcd}[ampersand replacement=\&]
e \arrow[r,dash] \& K \arrow[r,dash] \& \unsure{G|C} \arrow[r,dash] \& G' \arrow[r,dash] \& e
\end{tikzcd}}
(\unsure{non} commutatifs), tels que \deleted{$K$} \inserted{$G,G'$} \uline{simples/$k$}, on a
\[
H^1(k,G)\longrightarrow H^1(k,G')\qquad\text{\uline{surjectif}}
\]
% slanted
\margin{Inutile de \unsure{supposer} $G,G'$ \unsure{de t.f.|simples}. On se \unsure{ramène} \unsure{d'abord} \uline{\unsure{au}} \uline{\unsure{cas}} \uline{$K$ \unsure{réduit}} \unsure{? Si oui}}

Principe de démonstration. On prend un tore maximal $T$ de $K$ défini sur $k$, et utilisant \unsure{le|les} th. de conjugaison, on voit que $N_G(T)\longrightarrow G'$ est \deleted{surjectif} \inserted{\unsure{épim.}}, donc on aura
% slanted
\margin{faut-il supposer $K$ linéaire??? On \unsure{peut} commencer par \unsure{dévisser} en \uline{groupe fini}, et \deleted{groupe} \inserted{\unsure{groupes}} \uline{commutatifs} \uline{connexes} (\unsure{ou commutables}) \uline{et} groupe \uline{\unsure{quasi-réductif}} \unsure{Opérer ici leur propre} \unsure{comm.} \unsure{aux} \uline{\unsure{linéaires}}}
\diagram{\begin{tikzcd}[ampersand replacement=\&]
e \arrow[r,dash] \& N_K(T) \arrow[r,dash] \& N_G(T) \arrow[r,dash] \& G' \arrow[r,dash] \& e
\end{tikzcd}}
donc on peut supposer \uline{$T$ invariant dans $G$}, a fortiori dans $K$. Car on \unsure{se} ramène aux deux cas

1) $K$ est un tore \deleted{\unsure{1)}} $K$ est \deleted{sans \unsure{tore}} \inserted{de \unsure{rang}} \unsure{réductif} 0. \unsure{Par} dévissage ce dernier cas se réduit à :

2) $K$ est fini ét. /$k$ \deleted{\unsure{3)}} 3) $K$ est connexe, et géométriquement une ext. d'une VA par un unipotent. Par dévissage avec les commutateurs successifs, on est ramené
% archive sheet number: 3

